% ---------------------------------------------------------------------
%  CHƯƠNG 5 — KẾT QUẢ VÀ ĐÁNH GIÁ
% ---------------------------------------------------------------------
\chapter{Kết quả và đánh giá}
\label{ch:ketqua}

\begin{muctieu}
Chương này đối chiếu kết quả đạt được với từng mục tiêu đã đặt ra ở
Chương~\ref{ch:modau}, trình bày các số liệu đo lường định lượng từ bộ kiểm
thử tự động, phân tích sâu các kịch bản tranh chấp đa luồng và nêu các hạn
chế hiện tại của hệ thống.
\end{muctieu}

\section{Kết quả đạt được}

Đối chiếu trực tiếp với năm mục tiêu đã đề ra tại Chương~\ref{ch:modau}, đồ án
đã hoàn thành trọn vẹn tất cả các hạng mục:
\begin{enumerate}[leftmargin=1.6em, itemsep=2pt]
  \item \textbf{Đạt Mục tiêu 1 (Máy chủ TCP bể luồng):} \code{GameServer} sử
        dụng bể luồng cố định 20 luồng phục vụ ổn định các kết nối song song,
        tự động xếp hàng kết nối thứ 21 và thu hồi toàn bộ tài nguyên khi gọi
        \code{close()}.
  \item \textbf{Đạt Mục tiêu 2 (Giao thức văn bản đầy đủ):} Toàn bộ 12 lệnh
        từ phía máy khách và các bản tin \code{INFO}, \code{OK}, \code{ERROR}
        từ máy chủ hoạt động chính xác, xử lý an toàn mọi dòng rỗng, lệnh lạ
        hoặc tham số sai định dạng số.
  \item \textbf{Đạt Mục tiêu 3 (Hoàn thiện 21/21 chức năng):} Đầy đủ 21 chức
        năng từ \code{F01} đến \code{F21} đã được cài đặt và hợp nhất lên
        nhánh \code{main}.
  \item \textbf{Đạt Mục tiêu 4 (Giao diện \code{ClientUI} mượt mà):} Giao diện
        Swing tách luồng đọc mạng riêng, cập nhật trạng thái tức thời và hỗ
        trợ mở ván, đoán số, tra cứu bảng xếp hạng, lịch sử, thống kê ngay
        trên cửa sổ.
  \item \textbf{Đạt Mục tiêu 5 (Kiểm thử TDD 100\%):} Toàn bộ 103/103 ca kiểm
        thử tự động JUnit 5 đều vượt qua với 0 lỗi.
\end{enumerate}

\section{Độ đo đánh giá định lượng}

Bảng~\ref{tab:ketqua} tổng hợp các chỉ số định lượng thực tế thu được khi biên
dịch và chạy bộ kiểm thử tự động trên nhánh \code{main} của dự án.

\begin{table}[H]
\centering
\caption{Các chỉ số đánh giá định lượng của hệ thống trên bộ kiểm thử thực tế.}
\label{tab:ketqua}
\begin{tabular*}{\textwidth}{@{\extracolsep{\fill}}>{\bfseries\color{daunavy}}l c c l}
\toprule
\rowcolor{gray100}\bfseries\color{daunavy}Chỉ số đánh giá & \bfseries\color{daumaroon}Mục tiêu & \bfseries\color{daumaroon}Thực đạt & \bfseries\color{daunavy}Phương pháp kiểm chứng \\
\midrule
Số chức năng hoàn thành (\code{F01}--\code{F21}) & 21 / 21   & 21 / 21 (100\%)   & Đối chiếu bảng đặc tả Đề tài 30 \\
Số ca kiểm thử JUnit 5 vượt qua                  & 100\%     & 103 / 103 (100\%) & Chạy \code{junit-platform-console-standalone} \\
Thời gian thực thi toàn bộ 103 tests             & $< 10$~s  & $4,89$--$5,88$~s  & Đo trực tiếp bằng JUnit Platform \\
Độ chính xác chốt người thắng (100 luồng)        & Tuyệt đối & 1 người thắng duy nhất & Lặp nhiều vòng trong \code{GameDongThoiTest} \\
Kiểm thử chống \code{Deadlock} chuyển phòng      & 0 bế tắc  & 2.000 lượt / 40 luồng & \code{QuanLyPhongVaDaPhongTest} (timeout 5~s) \\
\bottomrule
\end{tabular*}
\end{table}

\section{Bàn luận}

Kết quả thực nghiệm cho thấy ba quyết định thiết kế mang lại hiệu quả rõ rệt
nhất cho độ ổn định của hệ thống:
\begin{itemize}[leftmargin=1.6em, itemsep=2pt]
  \item \textbf{Đồng bộ hóa ở mức hạt mịn \en{(fine-grained synchronization)}:}
        Thay vì đặt một khóa toàn cục trên toàn bộ máy chủ, mỗi đối tượng
        \code{Game} và \code{PhongChoi} tự quản lý khóa riêng của mình, đồng
        thời việc phát tin ra từng kênh \code{PrintWriter} chỉ khóa trên chính
        đối tượng \code{PrintWriter} của khách đó. Cách làm này giúp các phòng
        chơi khác nhau hoạt động song song thực sự mà không chặn lẫn nhau.
  \item \textbf{Giải quyết tranh chấp giữa hẹn giờ hết ván và lượt đoán trúng
        (\code{F20}):} Khi tác vụ hẹn giờ \code{Callable<Boolean>} thức dậy
        sau \code{Thread.sleep} đúng vào khoảnh khắc một người chơi gửi số
        đoán trúng, cả hai luồng đều phải đi qua cùng một khóa
        \code{synchronized (khoaVan)} và kiểm tra trạng thái \code{PLAYING}
        của \code{Game}. Luồng nào vào khóa trước sẽ chuyển trạng thái sang
        \code{FINISHED}; luồng vào sau thấy trạng thái đã là \code{FINISHED}
        nên tự động rút lui, không bao giờ xảy ra hiện tượng vừa công bố người
        thắng vừa báo hết giờ.
  \item \textbf{Xoay vòng lượt đoán bằng \code{ArrayDeque} (\code{F19}):}
        So với cách dùng danh sách kèm biến chỉ số \code{currentIndex}, việc
        giữ người đến lượt luôn ở đầu hàng đợi (\code{peek()}) và xoay vòng
        bằng cặp thao tác \code{poll()}/\code{offer()} giúp xử lý cực kỳ gọn
        khi người đang giữ lượt rớt mạng (\code{remove(ten)} tự động đẩy người
        kế tiếp lên đầu hàng) hoặc khi có người mới tham gia giữa ván
        (\code{offer(ten)} xếp vào cuối hàng).
\end{itemize}

\section{Hạn chế}

\begin{luuy}
Bên cạnh những kết quả đạt được, hệ thống hiện còn hai hạn chế thuộc phạm vi
giới hạn của môn học:
\begin{enumerate}[leftmargin=1.5em, itemsep=2pt]
  \item \textbf{Lưu trữ dữ liệu trong bộ nhớ \en{(in-memory)}:} Bảng điểm
        \code{KhoBangDiem} và lịch sử \code{KhoLichSuDoan} được giữ trong RAM
        của tiến trình \code{GameServer}; khi tắt hẳn máy chủ thì dữ liệu sẽ
        được khởi tạo lại từ đầu.
  \item \textbf{Kênh truyền TCP văn bản thuần:} Dữ liệu truyền giữa
        \code{ClientUI} và \code{GameServer} chưa được mã hóa bằng TLS/SSL và
        chưa có bước xác thực mật khẩu tài khoản.
\end{enumerate}
\end{luuy}
